\documentclass{article}

% Projektový záměr pro Komisi – výzkumný projekt.
% Záměr: oblast, problém, stávající stav, týmový prvek, rozsah a období. Bez opakování téhož v několika sekcích.

\usepackage[a-2u]{pdfx}
\usepackage[utf8]{inputenc}
\usepackage[T1]{fontenc}
\usepackage{array}
\usepackage[ddmmyyyy]{datetime}
\usepackage{geometry}
\usepackage{hyperref}

\renewcommand{\dateseparator}{.}
\renewcommand{\baselinestretch}{1.5}

\geometry{margin=2cm}

\def\ResearchProject{Výzkumný projekt}
\def\CompanyProject{Firemní projekt}
\def\SoftwareProject{Softwarový projekt}

% METADATA SECTION START
\def\StudentsFullName{Pavol Rajczy}
\def\Obor{Computer Science -- Software and Data Engineering}
\def\ProjetTitle{Podpora návrhu ontologií s asistentem a integrace s nástrojem Dataspecer}
\def\ProjectType{\ResearchProject}
\def\SupervisorFullName{prof. Mgr. Martin Nečaský, Ph.D.}
\def\ConsultantsFullName{Mgr. Štěpán Stenchlák}
\def\ExpectedStart{2025-03-01}
\def\ExpectedEnd{2026-09-01}
% METADATA SECTION END

\title{\ProjetTitle}
\author{\StudentsFullName}

\begin{document}

\centerline{\Large \textbf{Návrh projektového záměru}}
\centerline{\textbf{Katedra softwarového inženýrství}}
\centerline{\textbf{Matematicko-fyzikální fakulta, Univerzita Karlova}}
\bigskip
{\noindent\begin{tabular*}{\textwidth}{ >{\raggedleft}m{4cm} l}
 {\bf Řešitel:} & \StudentsFullName \\
 {\bf Studijní program:} & \Obor \\
 & \\
 {\bf Název projektu:} & \ProjetTitle \\
 {\bf Typ projektu:} & \ProjectType \\
 & \\
 {\bf Supervizor:} & \SupervisorFullName \\
 {\bf Konzultant(i):} & \ConsultantsFullName \\
 & \\
 {\bf Předpokládaný začátek:} & \ExpectedStart \\
 {\bf Předpokládaný konec:} & \ExpectedEnd \\
\end{tabular*}}

\medskip
\noindent\textbf{Souhlas:} Se záměrem projektu souhlasí všichni uvedení (řešitel, supervizor, konzultant). Souhlas může být doložen např. uvedením do kopie e-mailu se záměrem nebo podpisem na vytištěném záměru.

\section{Úvod}

V oblasti sémantického modelování dat slouží \emph{datové specifikace} k popisu struktury a významu dat tak, aby je různé systémy mohly konzistentně interpretovat a validovat. Specifikace zahrnují slovníky, aplikační profily a z nich odvozené technické artefakty -- schémata pro XML, JSON, CSV a dokumentaci. Návrh takové specifikace, zejména návrh ontologie nebo slovníku, je náročný: vyžaduje identifikaci pojmů, vztahů a omezení z odborných nebo právních dokumentů a jejich formalizaci.

Pro usnadnění této práce vznikl ve výzkumné skupině vedoucího projektu systém pro \emph{agentické sémantické modelování}: využívá velké jazykové modely (LLM) k iterativnímu návrhu ontologie z textu -- např. z právních předpisů uložených v \emph{knowledge base}. Knowledge base zde znamená soubor dokumentů (právní texty, odborná literatura), ze kterých asistent navrhuje třídy, atributy a vztahy; uživatel pracuje v tzv. design projektu, kde má k dispozici navržené doménové oblasti, iterace a konkrétní operace nad ontologií. Předchozí vývoj ukázal, že takový asistent je použitelný, ale zůstávají omezení: výstupy z asistenta nelze přímo využít v nástrojích, které výzkumná skupina a partneři používají pro správu specifikací; chybí ucelený workflow od návrhu po schválení a předání dál; chybí systematická zpětná vazba od uživatele, kterou by se asistent řídil.

\paragraph{Dataspecer.}
Dataspecer je nástroj pro správu sémantických datových specifikací: od slovníků a aplikačních profilů až po odvozené technické artefakty -- datová schémata, validační pravidla, API a prototypy aplikací. Umožňuje definovat třídy, vlastnosti a vztahy ve stylu RDF Schema, kombinovat termíny z existujících slovníků do aplikačních profilů a z doménové ontologie poloatomaticky odvozovat schémata pro XML, JSON a CSV -- tedy podporuje jak sémantickou, tak technickou interoperabilitu dat.

Cílem tohoto projektu je propojit asistenta s Dataspecerem a doplnit chybějící workflow: aby uživatel mohl návrhy z asistenta \emph{prohlížet, upravovat s využitím zpětné vazby (human-in-the-loop), schválit a výslednou ontologii exportovat} do Dataspeceru -- tedy aby celý proces od dokumentů v knowledge base až po použitelnou specifikaci v Dataspeceru byl možný bez ručního přepisování. Knowledge base by přitom neměla být pevně vázaná jen na asistenta; úložiště dokumentů může být např. jinde (včetně budoucí integrace s Dataspecerem), projekt na ni pouze odkazuje.

\subsection{Vztah k jiným projektům a týmový prvek}

Nástroj Dataspecer je projektem výzkumné skupiny. Pozostáva z rôzných nezavislých komponent spravonavých roznymi členmi týmu. Integrace asistenta s Dataspecerem a návrh workflow se dotýká backendu Daraspeceru, napríklad potrebou importovať a exportovať dáta z asistenta. Taktiež i jeho uživatelského rozhraní, aby sa užívateľ vedel dostať k asistentovi priamo z Dataspeceru, což vyžaduje spolupráci s členy skupiny zodpovědnými za Dataspecer.

Student se bude účastnit pravidelných  schůzek výzkumné skupiny a bude v součinnosti s členy skupiny navrhovat uživatelské scénáře, identifikovat požadované funkcionality a následně je implementovat, dokumentovat a testovat u uživatelů.

\section{Popis projektu}

\paragraph{Stávající stav, na který projekt navazuje.}
Existující agentický backend (Python, FastAPI) umožňuje: načtení dokumentů do knowledge base, identifikaci doménových oblastí, návrh iterací a plánů úloh a generování operací nad ontologií (přidání tříd, atributů, vztahů). Částečně existuje aplikace operací a perzistence ontologie. Backend nemá podporu výměny s Dataspecerem, chybí pracovní kopie ontologie (uživatel nemůže bezpečně experimentovat bez rizika přepsání původní verze), chybí náhled dopadu operací před schválením, mechanismus human-in-the-loop (zbírka oprav a preferencí uživatele, kterými se asistent řídí, včetně možnosti záznamy prohlížet a odstraňovat) a volitelně zacílení asistenta na část ontologie nebo zadání cíle. \emph{Podoblast ontologie} znamená vybranou část modelu (např. jen třídy a vztahy odpovídající jedné kapitole zákona nebo jednomu doménovému tématu), na kterou asistent omezí své návrhy -- uživatel pak dostává návrhy jen pro tuto část, ne pro celou ontologii. \emph{Cíl} je uživatelský požadavek typu „doplň vztahy mezi třídami X a Y“ nebo „vytvoř třídu Z s atributy z oddílu 3 zákona“.

\paragraph{Cíl projektu a co bude uživatel umět.}
Rozšířit systém tak, aby byl asistent použitelný v celém workflow až po export do Dataspeceru. Konkrétně:
\begin{itemize}
 \item Uživatel bude moci pracovat nad \emph{kopií} ontologie a experimentovat s návrhy bez rizika poškození původní specifikace.
 \item Uvidí \emph{náhled změn} před tím, než je potvrdí, a před exportem do Dataspeceru.
 \item Zpětná vazba (opravy, schválení/odmítnutí, preference) bude uložena jako \emph{záznamy, kterými se asistent řídí}; uživatel je bude moci prohlížet a odstraňovat položky, které už nechce.
 \item Výslednou ontologii bude možné \emph{exportovat} do formátu importovatelného v Dataspeceru.
 \item Volitelně: zacílení na podoblast a zadání cíle, aby asistent navrhoval jen ve vybrané části nebo směřoval k zadanému cíli.
\end{itemize}

\paragraph{Rozsah a scope.}
Rozsah zahrnuje: návrh vazby na knowledge base (ne pevné vázání na asistenta; úložiště může být např. v Dataspeceru nebo jinde) a vazeb projekt--knowledge base--ontologie; zavedení pracovní kopie a náhledu změn (preview/diff); human-in-the-loop -- návrh metod sběru zpětné vazby, reprezentace záznamů pro asistenta, jejich zobrazení a odstranění v UI a integrace do kontextu agentů; export do formátu pro Dataspecer a workflow schválení $\to$ export; volitelně podoblast a cíle; dokumentaci a testování s uživateli. Použité technologie: stávající backend (Python, FastAPI), ontology store, knowledge base, formát Simplified Semantic Model; rozšíření API a uživatelského rozhraní.

\end{document}
